\documentclass[10pt,conference]{IEEEtran}
\usepackage{cite}
\usepackage{amsmath}
\usepackage[T1]{fontenc}
\usepackage{url}
\usepackage{hyperref}
\usepackage{booktabs}
\usepackage{tabularx}
\usepackage{array}
% Author details preserved from paper.tex; excluded from anonymous rendering:
% Vitor Raposo, Ontario Tech University, Oshawa, Ontario, Canada
% vitor.raposo@ontariotechu.net
% Sanaa Alwidian, Ontario Tech University, Oshawa, Ontario, Canada
% sanaa.alwidian@ontariotechu.ca
% Study repository:
% https://github.com/MyLittleFoxxie/MESA-meta-evaluation-standard-for-ai

\begin{document}
\title{MESA: Adapting the EFPA Test Review Model for AI Benchmark Review}
\author{Anonymous authors}
\maketitle
\begin{abstract}
\textbf{Background:} AI benchmark scores alone do not establish whether a benchmark measures what it claims to. Psychological and educational testing offers established methods for reviewing tests and interpreting their scores.
\textbf{Objective:} We adapt the 90-page European Federation of Psychologists' Associations Test Review Model (EFPA-TRM), developed for psychological and educational tests, into MESA, a 25-question questionnaire intended to make AI benchmark review more accessible. We examine which criteria carry over and what detail is lost.
\textbf{Method:} Claude Sonnet 5.5 assesses information about ARC-AGI-2's grid puzzles and DesignQA's engineering questions using MESA, EFPA-TRM, or general review instructions in 30 separate conversations. Constructed test cases ask whether information confirms a claim, contradicts it, or leaves it unknown; expected answers are fixed before testing. Selected benchmark reviews are examined for information they identify, omit, or interpret too broadly.
\textbf{Results:} 28 of 30 conversations met the required format. Test-case answers matched the expected answers for 243 of 245 with MESA, 237 of 245 with EFPA-TRM and 191 of 196 with general instructions; every mismatch called a claim contradicted where the description left it unknown. MESA choices varied for 9 of 50 question--benchmark pairs, often between No and Not sure. Selected reviews under every method included conclusions needing qualification.
\textbf{Conclusion:} MESA expresses selected EFPA-TRM checks as 25 explained questions. The study does not show that it identifies more problems than other methods or that people find it easier to use; human evaluation is future work.
\end{abstract}
\begin{IEEEkeywords}
AI benchmarks, test review, questionnaire, software engineering, EFPA-TRM, score interpretation
\end{IEEEkeywords}

\section{Introduction}
A test score might seem self-explanatory, but it needs an explanation of what it tells us. For example, solving random selected puzzles may not by itself show that an AI can solve problems in other settings. Psychology and education have developed methods for examining how tests are built, how answers are marked, whether scores are consistent, and what their results mean \cite{aera_apa_ncme_standards_2014,kane_validation}. These methods provide a starting point for reviewing AI benchmarks themselves. We ask which established questions can be used for AI, which need a different meaning, and what a shorter questionnaire leaves out.

Human tests and AI benchmarks both use selected tasks and scoring rules to support conclusions, but their subjects and conditions differ. AI systems can encounter test material during training, use tools or supplied documents, and run under changing software settings. Benchmark critiques identify limits in task coverage, comparison conditions, and broad claims made from scores \cite{jo_wilson_what_does_your_benchmark_measure,singh_leaderboard_illusion}.

The European Federation of Psychologists' Associations Test Review Model (EFPA-TRM) provides the source for this adaptation. It separates a description of a psychological or educational test from an evaluation of the evidence supporting its use. The original 2025 document spans 90 pages, including forms and guidance \cite{efpa_test_review_model_2025}. Its established history motivates this investigation; it does not prove that the guidance works unchanged for AI.

Existing research provides guidance for assessing AI benchmark quality and interpreting benchmark scores \cite{reuel_betterbench,bean_measuring_what_matters,jo_wilson_what_does_your_benchmark_measure}. This paper examines a specific adaptation question: how can EFPA-TRM, developed for reviewing psychological and educational tests, be adapted to a test review model for AI benchmarks? We present MESA and explain which checks are retained, changed, combined, omitted, or added.

MESA, the Meta-Evaluation Standard for AI Benchmarks, is a 25-question questionnaire adapted from the EFPA-TRM. Each question has five answer choices, an explanation shown beside it, and an optional comment. No total score is calculated. The accompanying survey provides benchmark descriptions, explanations of terms, and a completed example. This format is intended to make review more accessible to people who are not specialists in test review. Here, Claude Sonnet 5.5 generates reviews from the supplied descriptions. The experiment does not run the benchmarks themselves, and human evaluation remains future work.

One possible application is a team examining whether benchmark results support selecting an AI component, a setting in which research on requirements engineering for AI reports continuing difficulties in specification and validation \cite{habiba_re4ai_maturity_2024}. Effects on team decisions are not measured here.

The paper contributes three deliverables: (1) MESA, a 25-question questionnaire for AI benchmark review; (2) an explanation of the adaptation decisions, including changes in meaning and the detailed checks combined or left out; and (3) an initial model-based examination showing examples of what the resulting reviews retain and leave unresolved. The explanation of the adaptation and its trade-offs is the central research contribution. The questionnaire makes those decisions usable in a concise format intended to make review easier to approach. The research questions are:
\begin{enumerate}
\item \textbf{RQ1:} To what extent can the EFPA-TRM be adapted into a concise questionnaire for reviewing AI benchmarks (MESA)?
\item \textbf{RQ2:} Do MESA-guided reviews distinguish between supported claims, documented problems, and missing information?
\item \textbf{RQ3:} What is retained or lost when AI benchmarks are reviewed with the MESA questionnaire rather than the full EFPA-TRM?
\end{enumerate}
In these questions, a supported claim is a conclusion justified by the supplied benchmark description. A documented problem is a problem explicitly described there, rather than a possible risk or information that the description does not provide. We examine these distinctions by giving the same benchmark descriptions to Claude Sonnet 5.5 with three sets of review instructions: use MESA, use the original EFPA-TRM, or follow general instructions to examine benchmark quality. We also use 24 constructed test cases to check particular distinctions. The study does not establish improved review accuracy.

\section{Background and Related Work}
\subsection{Learning from human-test review}
The paper draws on established methods for reviewing psychological and educational tests. The federation dates the first version of its review model to 2002, drawing on earlier national review systems \cite{efpa_test_review_model_2025}. EFPA-TRM first describes a test and its materials, then examines its purpose, development, score comparisons, consistency, support for its intended use, fairness, and reporting \cite{efpa_test_review_model_2025}. In its terminology, reliability concerns consistency and uncertainty in scores, while validity concerns the evidence supporting what scores are taken to mean. EFPA-TRM provides detailed checks and numerical ratings where appropriate. Its descriptions and evaluations serve different purposes: knowing how a test is given does not show that its results are consistent.

Applying this guidance to AI requires decisions about what the original questions mean. A score compared with results from a defined group of people does not directly become a ranking of AI systems. Comparisons need named systems and stated testing conditions. A report about a person differs from a public leaderboard, and expertise in using a human test differs from expertise in running and interpreting evaluation software. The EFPA-TRM comparison records these changes of meaning rather than automatically marking human-specific questions as failures.

\subsection{Score meanings and AI benchmark review}
A score and the claim made from it are different things. Evidence is needed to show that a score supports a particular meaning and use \cite{kane_validation,aera_apa_ncme_standards_2014}. In AI benchmarking, Bean et al. examine how the ability of interest, selected tasks, scoring rules, and resulting claims relate to one another \cite{bean_measuring_what_matters}. A review should ask both whether scoring follows its rules and whether the result supports the claimed ability. Correctly applying a scoring rule alone cannot answer the second question.

Benchmark critiques identify narrow task coverage and claims that go beyond the reported results \cite{line_goes_up,alonso_church_sigir}. Broad intelligence claims need evidence across the abilities included in the claim. Work drawing on Cattell--Horn--Carroll theory, an account of human cognitive abilities, provides one way to examine that breadth \cite{hendrycks_definition_agi}. A narrow benchmark can still be useful when the proposed claim is equally narrow.

BetterBench already provides a framework and checklist for assessing AI benchmark quality, covering design, implementation, documentation, maintenance, and interpretation of results \cite{reuel_betterbench}. MESA overlaps with these concerns and draws on BetterBench for software and maintenance questions. Bean et al. already examine whether benchmark tasks and scores justify the abilities claimed, and provide practical recommendations \cite{bean_measuring_what_matters}. Jo and Wilson likewise draw on measurement ideas to examine what benchmark results establish about AI capabilities \cite{jo_wilson_what_does_your_benchmark_measure}. Datasheets and Model Cards provide further ways to document datasets and models \cite{gebru_datasheets_2021,mitchell_model_cards_2019}.

These studies establish that benchmark review, practical checklists, and learning from human measurement research are existing lines of work. MESA's specific focus is the adaptation of EFPA-TRM into a shorter questionnaire and the consequences of those adaptation decisions. The comparison with the original model examines the detail represented in the questions and the extra explanation needed when applying them. This is a study of a particular adaptation, rather than evidence that no earlier work has adapted human-test review methods to AI.

\subsection{Models as reviewers}
Language models have been studied as judges of generated answers. Zheng et al. identify biases and compare model decisions with human preferences in their evaluation setting \cite{zheng_llm_judge_2023}. Shi et al. distinguish whether decisions remain the same across repeated evaluations from whether they change when answers appear in a different order \cite{shi_judging_judges_2025}. These findings motivate repeating the reviews, but do not establish that a model can review what a benchmark measures and what its scores justify. Here the model examines descriptions of benchmarks, rather than marking answers produced by the systems tested on them.

\section{The MESA Questionnaire}
\subsection{How EFPA-TRM concerns become questions}
Table~\ref{tab:domains} links the questionnaire's six areas to concerns in EFPA-TRM. A supporting document explains the proposed link for each question. These links explain the current design; they do not reconstruct an undocumented development history or imply endorsement by the federation. RQ1 concerns which checks can be expressed, how they change for AI benchmarks, and the difficulties observed when using them.

\begin{table*}[t]
\caption{MESA question groups and their proposed links to EFPA-TRM. Numbers identify questionnaire questions.}
\label{tab:domains}
\small
\begin{tabularx}{\textwidth}{@{}p{0.19\textwidth}p{0.07\textwidth}XX@{}}
\toprule
Question group & Questions & What EFPA-TRM examines & Change for AI benchmarks \\
\midrule
Description & 1--5 & Which test is being reviewed, its tasks, conditions, scoring, reports, and documentation & Ask whether key information is documented; passages supply the facts. \\
Purpose and development & 6--9 & Purpose, development, problem quality, coverage of the claimed ability & Keep reasons for selecting and checking tasks; combine questions about how well they cover the ability. \\
Scoring and comparisons & 10--13 & Scoring, consistent marking, combined scores, comparison groups & Simplify marking checks and ask whether comparison points fit the claim. \\
Consistency of results & 14--16 & Whether scores are consistent and how precise they are & Ask about repeated runs, changed conditions, and reported variation. \\
Meaning and fairness & 17--21 & What scores mean and whether comparisons are fair & Keep limits on claims; add attention to earlier exposure to questions and training on the benchmark. \\
Use and reporting & 22--25 & Materials, access, documentation, reports & Combine access and reuse concerns; ask what can be inspected, how changes are recorded, and what public claims say. \\
\bottomrule
\end{tabularx}
\end{table*}

The adaptation keeps questions about purpose, task selection, scoring, comparison points, and the limits of claims. It combines related details when one question can draw attention to a wider concern. For example, a question about recorded changes and their effects on comparisons covers several kinds of benchmark updates. Responses use five choices with optional comments.

Some detail is omitted. MESA does not separately assess the qualifications of people using human tests, tables comparing a person with a defined group, or methods examining how individual questions relate to each other and to the ability being measured. These methods may still matter for AI. Asking how scores change across runs does not answer how well the individual questions work together. The comparison with EFPA-TRM makes such omissions visible.

Questions added for AI concern earlier exposure to test questions, training specifically for a benchmark, inspection of software, and changes in evaluation procedures. BetterBench informs software and maintenance concerns \cite{reuel_betterbench}; measurement research informs whether tasks and scores support the claims \cite{bean_measuring_what_matters}. These additions do not imply that every AI issue has a matching original EFPA-TRM question.

\subsection{Three adaptation choices}
\textbf{Checking provided answers and marking (questions 8 and 12).} Human-test review separates checks of test content from checks of consistent marking \cite{efpa_test_review_model_2025}. For AI benchmarks, provided answers can contain mistakes, while software or another AI may mark new responses. Software can follow its scoring rule even when a provided answer is wrong. MESA question 8 asks whether problems and the answers used to mark them were checked for mistakes. Question 12 asks whether marking was tested for correct and consistent scores. This keeps the concerns separate without recording every kind of check, how often reviewers agreed, or how many mistakes remained.

\textbf{Choosing comparison points (question 13).} Human-test review examines how a person's score is compared with results from a defined group and whether that group fits the intended use. For AI, comparing available systems does not establish that they represent all systems of interest. MESA question 13 asks whether a human score, another AI's score, or a guessing result is suitable for the claim. It leaves out detailed requirements for selecting a wider comparison group or setting a score cutoff used to decide whether a requirement is met. Those decisions need further information; a leaderboard alone does not establish either.

\textbf{Repeating a test and changing its conditions (questions 14--16).} Human-test review separates repeated testing from comparisons between test versions intended to give comparable results. For AI, repeated runs, changed prompts, supplied documents, and software settings are different possible sources of score changes. MESA question 14 concerns repeated runs under the same conditions, question 15 concerns changed instructions or setup, and question 16 concerns whether estimates of variation are detailed enough for the comparison. A changed prompt does not establish that test versions measure the same thing in the same way. An uncertainty range calculated from the number of questions also need not describe changes across repeated runs. Detailed methods for studying how individual questions work together remain outside these short questions.

These examples make the adaptation decisions concrete. Checking the provided answers remains separate from checking marking; comparing scores retains a broad question but leaves out detailed requirements for a wider comparison group or a decision cutoff; and repeated evaluations remain distinct from changes in testing conditions. Each choice determines what the short questionnaire asks explicitly and what needs further examination. The current design has 25 questions; the study does not establish that this is the best length or that it retains every distinction needed for a decision.

\subsection{Accessibility as a design goal}
MESA is intended to make benchmark review easier to approach through 25 questions and familiar answer choices. The aim is to reduce the reading and interpretation required to use the complete 90-page EFPA-TRM, including its forms and guidance, so that people who build, select, or question AI benchmarks can review them without first mastering a specialist instrument. Question count and document length describe the materials; they do not establish lower human effort or easier use, which the model study does not measure.

Accessibility shapes four design choices. Each question is shown with its explanation beside it rather than in a separate manual. Answers use five plain categories instead of numerical ratings. Comments are optional, so a review does not require written evidence for every answer. A completed example demonstrates the questionnaire on another benchmark. Each explanation states what the question asks and what an answer does not establish. For example, evidence that marking was tested does not establish complete accuracy, and missing results from repeated evaluations do not demonstrate unstable scores. Optional prompts invite reviewers choosing Partly to identify which part is supported, those choosing No to name information showing the requirement was not met, and those choosing Not sure to name the missing information. Three closing reminders distinguish missing information from poor quality, testing that occurred from thorough testing, and task results from broader claims.

\subsection{Questions, choices, and reading support}
The six question groups contain five, four, four, three, five, and four questions respectively. Questions 1--5 ask whether the benchmark, the ability it claims to test, its intended use, tasks, testing conditions, and scoring and reporting are documented. A reviewer answers whether this information is provided; the benchmark description supplies the facts. Questions 6--25 ask why tasks were chosen, what was checked, and what the results justify.

All questions use five choices: \emph{Yes}, the information supports the whole condition; \emph{Partly}, it supports some parts; \emph{No}, it shows the condition is not met; \emph{Not sure}, information is missing or insufficient; and \emph{Not applicable}, the issue does not arise for the benchmark or intended use. An omitted response is different from all five. If a question asks whether testing took place, Yes means that it occurred even when its findings were unfavorable.

The choices are not numerical quality ratings and are not added or averaged into a score. Not sure and Not applicable do not mean low or high quality: absence from the supplied material does not prove that a study was never conducted.

The survey packet (version 2.0) presents an introduction, three familiarity self-reports, a completed Humanity's Last Exam (HLE) example with all 25 choices and short reasons, ARC-AGI-2, DesignQA, and submission. The example is not an answer key for other benchmarks. The model study uses the reading content but not the familiarity questions or the browser interface.

\section{Study Method}
\subsection{Research method and empirical standards}
The study examines the proposed adaptation by asking Claude Sonnet 5.5 to produce reviews and comparing its answers. This is engineering research: it proposes a method and examines its use. We use the ACM SIGSOFT General and Engineering Research standards, together with the Benchmarking standard for the repeated model evaluations \cite{sigsoft_empirical_standards}. The model reads supplied benchmark descriptions and produces reviews; it does not execute ARC-AGI-2 or DesignQA. The two benchmarks provide contrasting tasks, and the constructed test cases check whether the model distinguishes what a description reports, contradicts, or leaves unknown. The experiment compares the complete supplied materials for each method. Each setup was repeated five times to describe variation, and every output is retained.

\subsection{Benchmark descriptions and three review methods}
ARC-AGI-2 asks a system to infer a transformation from examples of coloured grids and apply it to new grids \cite{chollet_arc_agi_2_2025}. DesignQA asks questions about engineering rules and design documents \cite{doris_designqa_2024}. The supplied descriptions identify the benchmark editions, evaluation software, source inspection dates, and what the scores are intended to show. A completed Humanity's Last Exam (HLE) review demonstrates how to answer MESA's questions \cite{phan_etal_hle_2025}.

The experiment uses three review methods. Each receives identical benchmark descriptions and intended uses. With \emph{MESA}, the model receives the 25 questions, answer-choice definitions, explanations, and completed HLE example. With \emph{EFPA-TRM}, it receives the original EFPA-TRM: text extracted from all 90 pages of the 2025 PDF, with page markers. With \emph{general instructions}, called ordinary review in the records, it receives instructions to examine benchmark construction, scoring, comparisons, consistency, interpretation, fairness, and reporting. These instructions also ask it to distinguish problems explicitly described in the material from missing information.

Each review includes the output requested by its assigned method and a common short conclusion. That conclusion lists claims justified by the description, problems explicitly reported there, and information needed but not supplied. This common format allows the content of the reviews to be compared despite different questionnaire and rating formats. MESA requires an answer to each of its 25 questions. EFPA-TRM requests each numbered descriptive field and rated item; these may be grouped only when their numbers and a shared reason are given. The model must also identify EFPA-TRM questions that need a different meaning for AI or do not apply. Only MESA includes the completed HLE example, and the descriptions replace images with captions and written descriptions.

\subsection{Test cases for interpreting benchmark information}
Twenty-four constructed test cases provide short benchmark descriptions with 49 claims for the model to classify. For example, a claim asks whether a benchmark was tested repeatedly: a description can report repeated testing, explicitly report only one evaluation, or give no information about repetition. The model chooses \emph{Supported}, \emph{Contradicted}, or \emph{Not established}, meaning that the supplied description justifies the claim, contradicts it, or leaves it unknown. The model must quote the description and explain its answer using its assigned review method.

Cases C1--C5 ask about repeated evaluations and changes in scores. C1 provides no information about repetition; C2 explicitly reports one evaluation. C3 reports ten closely grouped scores under a stated rule for calling them stable; C4 reports ten widely varying scores under the same rule. Repeated evaluation therefore occurred in both C3 and C4, but only C3 meets the stated rule. C5 adds installation instructions to C1 without adding any evaluation results.

C6 and C7 give the same grid-task results but change the intended use from comparing performance on those tasks to making a broad claim about intelligence. The results still justify the narrow task comparison, but do not establish the broader claim. C8 describes initial test results followed by a randomized workplace trial examining whether decisions were correct. Its claims distinguish the improvement reported in that workshop from safety, which was not measured.

The other sixteen cases form eight pairs, P1--P8. The two versions of each pair differ in one detail and receive the same two claims. P1 contrasts no information about repetition with an explicit record of a single evaluation; P2 contrasts stable with widely varying repeated scores under a stated rule; P3 adds installation material that reports no further results. P4 reports that editors checked the provided answers, with and without a finding of wrong answers. P5 contrasts marking software that gives the same mark to identical text with marking tested against independently established marks. P6 changes the stated use from a narrow task comparison to a claim about general intelligence. P7 changes whether a whole-grid scoring rule matches the stated scoring requirement. P8 contrasts changed instructions with identical repeated evaluations. Six claims in the pairs state a negative conclusion that the description does not justify, such as treating an unreported repeat as evidence that scores are unstable.

The expected answer for each claim and the reason for it were specified before the recorded reviews and withheld from the model. The cases and expectations were written while developing the questionnaire and had been used in earlier model runs; they are not new held-out tests. Matching these expected answers tests the particular distinctions built into the cases. It does not establish that reviews of ARC-AGI-2 or DesignQA are correct.

\subsection{Model and separate conversations}
The study requested Claude Sonnet 5.5 (\texttt{claude-sonnet-5-5}) with medium effort through Claude Code's non-interactive command-line mode \cite{anthropic_claude_code_cli}. Each session reported this model identifier. Underlying model revisions beyond that identifier and the settings controlling random output choices were not exposed. The model received no assigned reviewer identity.

Each experimental execution starts a separate model conversation, called a session in the records, without previous review history. For each of the three methods, five conversations review ARC-AGI-2 followed by DesignQA, and five other conversations answer all 24 test cases. This gives 30 planned conversations: 15 about the published benchmarks with 30 benchmark reviews, and 15 about the test cases with 735 answers to 49 repeated claims. MESA accounts for ten benchmark reviews with 25 answers each: 250 planned choices.

Each conversation starts in an empty temporary directory. The available settings remove all tools, including file reading, command execution, and web access, and disable project instruction files, memories, plugins, extension servers, and start-up scripts. A one-sentence system prompt replaces the tool's default instructions and asks the model to use only the supplied text in the requested format; the supplied instructions also prohibit external reading. Each session's start record lists the tools it offers, and event logs record any tool use. These controls do not remove knowledge learned during model training.

\subsection{Records and analysis}
The protocol specifies the instructions, expected test-case answers, response format, case order, and handling of failures before execution. File lists and digital fingerprints identify each input; original responses, event logs, settings, errors, and completion status are retained. If an execution fails without a completed response, the protocol permits one retry in a new conversation and retains both attempts. Completed responses with an invalid format are not repaired or replaced.

For RQ1, we examine the proposed links between EFPA-TRM and MESA alongside difficulties reported in the generated reviews. For RQ2, we count how often the model's answers to the test-case claims match the answers specified before testing, separately for each method and conversation. We compare 23 pairs of related claims within each conversation: C1/C5 adds unrelated instructions, C3/C4 changes the spread of scores, C6/C7 changes the proposed claim, C8 contrasts initial results with a later trial, and P1--P8 each change one detail. We also count Supported answers to the six unjustified negative claims. For the published benchmarks, we count each MESA answer choice by question and benchmark. We also count how often choices match across the five repeated conversations. We do not add choices into quality scores or conduct statistical significance tests.

For RQ3, we examine the first conversation about the published benchmarks for each method. Each response reviews both benchmarks, giving six selected reviews. The analyst reads each complete review, including its conclusion lists, against twelve points about the benchmark descriptions listed in a comparison checklist fixed before execution. The examination asks whether a conclusion is justified by the supplied information, whether a checklist point is left unaddressed, whether a conclusion goes beyond the description, and whether an EFPA-TRM question needs a different meaning for AI. EFPA-TRM is a comparison instrument, not the correct-answer key. These six reviews illustrate differences; they do not establish how often relevant points went unaddressed. All usable MESA closing notes are also read for reported difficulties.

\section{Results}
\label{sec:results}
The study ran on 17 October 2026. All 30 planned conversations completed through 30 physical attempts, without failures or retries. 28 responses met the predefined format. Two completed responses were retained without repair or replacement and excluded from the structured counts. One EFPA-TRM benchmark review placed line breaks inside text fields, which the required format does not allow. One test-case conversation following general instructions added a case identifier that is not in the case set. Table~\ref{tab:usable} gives the denominators for each method.

\begin{table}[t]
\caption{Usable conversations and test-case answers matching the answers specified before testing. Each method had five planned conversations in each batch.}
\label{tab:usable}
\centering
\small
\begin{tabular}{@{}lrrr@{}}
\toprule
Method & Benchmark & Test-case & Expected answers \\
\midrule
MESA & 5/5 & 5/5 & 243/245 \\
EFPA-TRM & 4/5 & 5/5 & 237/245 \\
General instructions & 5/5 & 4/5 & 191/196 \\
\bottomrule
\end{tabular}
\end{table}

\subsection{Application of the adapted questions (RQ1)}
Table~\ref{tab:domains} and the examples above show which EFPA-TRM checks are represented in MESA and which require more detailed examination. In the selected DesignQA review, the MESA response considered checks of provided answers separately from tests of marking and listed reviewer disagreement rates and remaining errors as missing information. The selected EFPA-TRM review identified the same missing counts and also stated that no score cutoff had been set; the selected MESA review did not discuss cutoffs. Completing every answer does not establish that every relevant point was examined.

All five MESA conversations about published benchmarks produced both complete questionnaires: ten benchmark reviews and 250 choices. Their closing notes still reported uncertainty about some questions. Five of the ten notes named question 8, question 12, or both as hard to apply or to separate from related questions. Others reported overlaps between questions 5 and 10 and between 17 and 18, and judgments about how much evidence is enough for the human-comparison and fairness questions. Three DesignQA notes stated that missing information had been recorded as Not sure rather than No. RQ1 therefore yields an explained adaptation with examples of remaining decisions: how much checking a positive answer requires, and where one question ends and a related one begins.

\textbf{A recorded DesignQA example.} Questions and passages are quoted exactly; choices and comments are from the first MESA conversation and are not correct answers.

\begin{quote}
\textbf{Question 8:} Were the test problems and the answers used to mark them checked for errors during development?

\textbf{Passage:} ``The paper reports that two additional reviewers checked manually written questions, with exceptions for questions derived from others and some explanations.''

\textbf{Recorded choice: Partly.} Comment: ``Additional reviewers checked written questions with exceptions; errors remaining unknown.'' The answer recognizes that checking took place without treating it as evidence that every provided answer is correct.

\textbf{Question 12:} Was the marking process tested to see whether it awards the right scores consistently?

\textbf{Passage:} ``Small marking-test files and discussion of scoring problems show some checking, but do not establish a complete study of marking accuracy and consistency.''

\textbf{Recorded choice: Partly.} Comment: ``Small marking-test files and scoring discussion, but no full study.'' The answer treats limited marking tests as partial support for the question rather than as a test that took place. MESA's answer-choice definitions allow either reading, which is the interpretation problem reported in the closing notes.
\end{quote}

All five MESA conversations chose Partly for both questions on DesignQA.

\subsection{What the descriptions establish or leave unknown (RQ2)}
Fourteen usable conversations answered the test cases. They repeated the same 49 claims across 24 cases; the 686 answers are not independent test cases. Table~\ref{tab:controlled} reports matches with the answers specified before testing. MESA matched 243 of 245 answers, EFPA-TRM 237 of 245, and general instructions 191 of 196. Three MESA conversations, one EFPA-TRM conversation, and one general-instruction conversation matched all 49 claims.

\begin{table}[t]
\caption{Answers to the separate test-case claims. Counts describe repeated outputs of one model.}
\label{tab:controlled}
\centering
\begin{tabular}{@{}lrrr@{}}
\toprule
Review method & Conversations & Expected answers & All 49 matched \\
\midrule
MESA & 5 & 243/245 & 3/5 \\
EFPA-TRM & 5 & 237/245 & 1/5 \\
General instructions & 4 & 191/196 & 1/4 \\
\bottomrule
\end{tabular}
\end{table}

All 15 mismatches had the same form: the model answered Contradicted where the expected answer was Not established. Twelve concern two claims. P5A reports that no comparison of marks with independently established marks is reported; five answers (one MESA, one EFPA-TRM, three general instructions) called the claim that marking was tested contradicted. The description leaves testing unknown rather than showing that it did not occur; one MESA explanation noted this itself (``reads as not tested (strictly, not reported)''). P6B asks whether grid results are sufficient for a stated use about general intelligence; seven answers (one MESA, four EFPA-TRM, two general instructions) called that claim contradicted. Here the expected label is debatable, because the description itself says that no other results are supplied. The remaining three mismatches come from one EFPA-TRM conversation, which also called the broad claims in C6 and C7 and the safety claim in C8 contradicted.

All 322 comparisons between related claims changed or stayed the same as expected; in 13, one of the two labels differed from the key in the way just described. Per method, MESA had 113 of 115 comparisons with both expected labels, EFPA-TRM 109 of 115, and general instructions 87 of 92. No conversation supported any of the six negative claims that the descriptions do not justify (0 of 84 opportunities). All 686 quotations were found in the supplied descriptions. For RQ2, the model handled the distinctions built into the cases under all three methods; MESA had the fewest mismatches, but the differences are small, come from one model, and concentrate on two claims, so they do not show that MESA distinguishes claims better.

Answers about the published benchmarks varied. Nine of the 50 question--benchmark pairs had different choices across the five MESA conversations (Table~\ref{tab:variation}). Of the 500 pairs of repeated answers (ten per question and benchmark), 456 (91.2\%) matched; these pairs are not independent. Four of the nine varying pairs moved between No and Not sure, the distinction between information showing a requirement is not met and information that is missing. Across all 250 choices, the model chose Yes 59 times, Partly 141, No 35, Not sure 15, and Not applicable never; matching choices do not establish correctness.

\begin{table}[t]
\caption{Question--benchmark pairs with varying MESA choices across five conversations.}
\label{tab:variation}
\centering
\small
\begin{tabular}{@{}llp{0.52\columnwidth}@{}}
\toprule
Benchmark & Question & Choices \\
\midrule
ARC-AGI-2 & 7 & Yes 3, Partly 2 \\
ARC-AGI-2 & 14 & No 4, Not sure 1 \\
ARC-AGI-2 & 15 & No 2, Not sure 3 \\
DesignQA & 6 & Yes 1, Partly 4 \\
DesignQA & 14 & Partly 3, No 2 \\
DesignQA & 16 & Partly 1, No 4 \\
DesignQA & 19 & No 1, Not sure 4 \\
DesignQA & 20 & No 3, Not sure 2 \\
DesignQA & 23 & Partly 1, No 4 \\
\bottomrule
\end{tabular}
\end{table}

\subsection{Comparison with EFPA-TRM (RQ3)}
The six selected reviews addressed most of the same points. All separated a completely correct output grid from partial credit across several inputs, performance on the benchmark tasks from broader abilities, and changed testing conditions from repeating the same evaluation. Against the twelve-point checklist, the selected EFPA-TRM and general-instruction reviews addressed all twelve points and the selected MESA review addressed ten and partly addressed two. The MESA review listed the matching of test sets for human difficulty and the missing repeat-run estimates separately, without stating that the matching leaves changes in AI scores unmeasured; the other two reviews stated this. It also listed DesignQA's always-yes possibility as a documented problem without noting that no tested AI was shown to use it; the general-instruction review added that ``no AI is shown to do this.''

Table~\ref{tab:review-tradeoffs} gives examples from these reviews. A detail absent from one response is not necessarily omitted by the questionnaire; by contrast, MESA does omit separate requirements for selecting a wider comparison group and setting a score cutoff.

\begin{table*}[t]
\caption{Selected MESA and EFPA-TRM outputs from the six selected reviews (first conversation for each method, both benchmarks). Differences describe these outputs, not the range of checks or accuracy of these approaches in general.}
\label{tab:review-tradeoffs}
\small
\begin{tabularx}{\textwidth}{@{}p{0.17\textwidth}XXX@{}}
\toprule
Review topic (MESA questions) & What the selected MESA review says & What the selected EFPA-TRM review adds & Why the difference could matter \\
\midrule
Provided answers and marking (8, 12) & DesignQA: checks of provided answers and marking tests are answered separately (both Partly); reviewer disagreement rates and remaining errors are listed as missing. & The question reviewers' unreported agreement is placed under the original model's agreement-between-markers item, while automated marking is treated separately. & Knowing that people checked the answers does not show how often they agreed or whether errors remain. \\
Scores used for comparison (13) & Random-guessing and AI comparisons are noted, without a human engineers' baseline. & No score cutoff was set; yes/no compliance answers rely on supplied answers rather than a decision standard. & Comparing particular systems does not establish a model's standing in a wider group or justify a decision cutoff. \\
Variation and changed conditions (14--16) & One AI tested with five passage sets gives some information about variation, but not about identical repeats. & Repeated runs and changed passage sets are mapped to the original reliability checks; one answer in the 16-question group changes accuracy by 6.25 percentage points. & In a small group of questions, one answer can substantially change the score; changing inputs also changes what is being compared. \\
Access and change records (23, 24) & No licence file, complete reuse rules, or clear change record were found; both questions are answered No. & The same findings are listed, adding that they do not show reuse is prohibited. & Information not found in the inspected materials is different from evidence that it does not exist. \\
\bottomrule
\end{tabularx}
\end{table*}

Using the original model required explicit changes from human testing to AI evaluation. The selected EFPA-TRM response treated the AI system with its settings, or its supply of information, as the subject taking the test. It treated the people who tried ARC-AGI-2's puzzles as a calibration of difficulty rather than a group for comparing scores, and marked most comparison-group requirements as not applicable or not possible to rate. It mapped repeated runs to consistency over repeated testing, and placed prior exposure and training on the test, which have no direct counterpart in the original model, under its validity and restriction checks. Seven of the eight usable EFPA-TRM closing notes asked whether a human-test question without an AI counterpart should receive a rating of 0 or be marked not applicable. MESA expresses several of these choices directly in its wording.

Each method's selected reviews included conclusions needing qualification. MESA answered No to DesignQA questions 23 and 24 because no licence file, reuse rules, or change record were found, although the passage states that these findings do not establish that permission is unavailable elsewhere. MESA's own instructions say not to convert missing information into No, and the same response's closing note says missing information was recorded as Not sure. For ARC question 14, it answered No from the single-run verified-testing policy, as did three other MESA conversations, while one chose Not sure; that policy does not rule out repeat-run studies by others. The review following general instructions listed permitted training on development problems, a stated policy, and the absence of repeat-run estimates as documented problems. The EFPA-TRM review listed the absence of studies predicting real-world performance as a documented problem while also calling those links unresolved. These examples distinguish a reported procedure, missing information, and evidence of a problem. They come from six selected reviews and do not estimate error rates.

For RQ3, the selected reviews identify mostly the same issues, differ in detail, and show how EFPA-TRM questions need reinterpretation for AI. No examined point appeared only in the MESA review, which left two points less explicit than the others.

\section{Discussion}
\subsection{What human-test review contributes to AI evaluation}
The central contribution is the explanation of the adaptation: which questions keep their meaning, which need a different meaning for AI, and which details are combined, omitted, or added. The experiment does not establish that the shorter questionnaire produces better reviews or covers the same checks.

The proposed adaptation and existing review examples suggest three lessons. First, checking the provided answers and checking how responses are marked are separate tasks. A mistake in either can affect the score, and checking one does not settle the other. Second, knowing that a check took place does not establish that it was thorough enough. A reviewer may need information about who checked what, how often they agreed, or what errors remained before relying on the finding. Third, comparing particular systems does not establish how one system stands among a wider group or justify a score cutoff for a decision. Such claims require more information about the comparison group and the reason for the cutoff.

These are lessons for using the questionnaire, not a tested procedure: a short question keeps a concern visible while leaving the depth of examination to the reviewer.

The study also shows where MESA's rules were not enough. Its explanations say not to convert missing information into No, yet four of the nine varying MESA pairs moved between No and Not sure, and a selected MESA review answered No where the passage reported only that information was not found. All test-case mismatches involved the same distinction: under all three methods, insufficient or unreported information was sometimes called contradicting evidence. A concise questionnaire can state the rule; it does not by itself ensure that the rule is applied.

The constructed cases check particular distinctions against fixed answers, while the published-benchmark reviews illustrate decisions about what is being reviewed, how much checking is enough, and what a broader claim needs. Completing all questions or repeatedly choosing the same answers does not establish that a review is correct.

\subsection{Possible uses in software-engineering work}
This approach could support benchmark development, examination of published claims, or discussions about an AI component. In the proposed review practice, an evaluator records an answer choice and explains which information justifies it. A colleague can then question whether the same evidence is sufficient for the team's proposed use. For a component expected to answer engineering-document questions, tested scoring software would still leave open whether the provided answers are correct and whether one rulebook represents local documents.

Comments are optional, so a completed questionnaire can lack the reasons a team needs to examine an answer. In the proposed practice, colleagues may need to request those reasons before relying on an answer for an important decision. A blank comment does not change the recorded choice, establish that the benchmark is poor, or justify inventing a reason. If reviewers choose Yes and Partly, their explanations may reveal different views about how much checking is enough or whether they are reviewing organizer-verified or community-reported evaluations. Effects on collaboration have not been observed, and repeated conversations with one model do not show how colleagues with different responsibilities would respond.

\subsection{Limitations}
The constructed test cases, expected answers, questionnaire text, and analysis scripts were prepared earlier with an AI assistant and had been used in earlier model runs. Claude Code with Claude Opus 5.5 conducted this study, examined the six selected reviews, and revised the manuscript. The reviewer model, Claude Sonnet 5.5, comes from the same model family as that assistant. Interpretations have not received independent human validation. The expected label for one claim (P6B) is debatable, and seven mismatches concern it. The study was not publicly registered before testing. The test-case checks do not validate all 25 questions or the quality of a complete review.

Coverage is restricted to two deliberately selected real benchmarks and short supplied passages. The six selected reviews illustrate differences without estimating how often they occur across all generated reviews. One completed EFPA-TRM benchmark review and one general-instruction test-case conversation were excluded because of their format; that does not establish that their conclusions were worse, and it leaves unequal denominators. Grouped EFPA-TRM item numbers do not prove that all checks were fully applied, and text extraction and omitted images may affect the comparison.

The methods receive different amounts of guidance, and only MESA receives the completed HLE example. Differences therefore cannot be attributed to questionnaire wording alone. The fixed order means that reading an earlier benchmark or test case may affect a later answer. Repeated reviews, answers to claims, and answer pairs come from one model and cannot be treated as independent people or models. Knowledge learned during training, model updates, and unavailable details about the model and random generation limit experimental control and exact repetition.

Generated comments describe questions the model reports as difficult to interpret. They do not measure human understanding or ease of use, or reveal why the model chose an answer. Some comments concern the experiment's instructions rather than MESA wording: two general-instruction notes pointed out that the DesignQA passage mentions a completed review that is not included, and one noted that the task mentions an assigned method when none is named. MESA's accessibility features were all supplied to the model as text, so the study cannot show whether they make review easier for people. Human review effort, understanding, and effects on decisions are outside the present findings.

\section{Artifacts and Future Evaluation}
The supporting materials include the protocol, fixed model instructions, benchmark descriptions, original-source details, expected test-case answers kept separate from reviewer inputs, execution records, and analysis outputs. A separate static HTML document presents the questionnaire with its explanations beside each question, without answer fields, saving, tracking, or submission. The live participant survey remains unchanged.

The repository guide \texttt{CHASE-study.md} links the protocol, exact inputs, raw responses, failures, and analysis records, and gives commands for reproducing the offline calculations: \url{https://github.com/MyLittleFoxxie/MESA-meta-evaluation-standard-for-ai}.
Future work includes clearer guidance for separating missing information from negative findings, test cases written independently of the questionnaire's authors, other models, additional benchmarks, and a human evaluation of perceived ease, interpretation difficulties, completion effort, and identification of important issues; faster completion alone would not establish a better review.

\section{Conclusion}
This paper examines how EFPA-TRM can be adapted to a more accessible test review model for AI benchmarks and what changes in that process. Its deliverables are MESA's 25 questions with explanations, the explanation of the adaptation decisions and their trade-offs, and an initial examination using model-generated reviews. The paper does not establish that MESA is more accurate or covers the same checks as EFPA-TRM.

The adaptation and selected reviews show why checking the provided answers should remain separate from checking how new responses are marked. They also distinguish evidence that a check took place from evidence that it was thorough enough, and a comparison of particular systems from a ranking in a wider group or a justified score cutoff. In the model study, all three methods matched most expected test-case answers; MESA matched slightly more, a difference too small and too dependent on one model to show an advantage. Every mismatch, and several MESA choices, treated missing or insufficient information as a negative finding, the distinction MESA's guidance is meant to protect. Accessibility remains a design goal rather than a measured human benefit; human evaluation is future work.

\section*{Use of AI}
OpenAI Codex was used to prepare earlier drafts of this manuscript, the model instructions, constructed test cases, expected answers, and analysis scripts \cite{openai_codex_execution}. Claude Code with Claude Opus 5.5 was used to conduct the recorded model conversations with Claude Sonnet 5.5, adapt the execution scripts, examine the selected reviews, and revise this manuscript \cite{anthropic_claude_code_cli}. The requested model identifier and available settings are recorded in the materials. The analysis and draft conclusions have not been independently checked by human reviewers. Neither assistant is an author, participant, or source of independent validation.

\bibliographystyle{IEEEtran}
\bibliography{references}
\end{document}
